\section{Panorama del curso y avances del aprendizaje profundo}

MIT 6.S191 es el curso introductorio oficial de MIT sobre aprendizaje profundo, impartido por Alexander Amini y Ava Amini. Se trata de un campamento intensivo (boot camp) de una semana de duración que cubre la teoría y la práctica fundamentales del aprendizaje profundo. 2025 es el octavo año en que se imparte el curso, que ha enseñado ya a más de 13 millones de estudiantes en todo el mundo.

\subsection{Una década de evolución del aprendizaje profundo}

El profesor Amini abrió la clase con una comparación elocuente: hace apenas una década (2015), el sistema de generación de rostros más avanzado en aprendizaje profundo solo producía bloques de píxeles borrosos (Goodfellow et al.); para 2018 (Karras, Laine, Aila), los rostros generados ya tenían un realismo fotográfico; en 2020, el equipo del curso MIT 6.S191 llegó incluso a producir un video deepfake en el que una figura de «Obama» daba la bienvenida a los estudiantes, video que acumuló más de 1.1 millones de reproducciones en YouTube.

\begin{figure}[H]
\centering
\includegraphics[width=0.85\textwidth]{slides-images/slide-002.jpg}
\caption{Una década de avances en la generación de rostros mediante aprendizaje profundo: de los píxeles borrosos de 2015 al video realista de 2020\protect\footnotemark}
\end{figure}
\footnotetext{Diapositiva 2.}

\begin{knowledgebox}{El costo de producción del deepfake de 2020}
Producir aquel video deepfake de 2 minutos en 2020 requirió: 2 horas de grabación de audio profesional, 50 horas de datos de video en alta definición, un guion estático predefinido y más de 15,000 dólares en recursos de cómputo. Para 2025, en cambio, Amini hizo una demostración en vivo de clonación de voz en tiempo real: bastaron unos minutos de audio de su clase grabados para generar de inmediato un clon de su voz, con el que sostuvo una conversación dinámica y sin guion. Esto ilustra con claridad el enorme salto que ha dado la IA generativa en tan solo unos años.
\end{knowledgebox}

\begin{figure}[H]
\centering
\includegraphics[width=0.85\textwidth]{slides-images/slide-006.jpg}
\caption{2020 contra 2025: comparación de las capacidades de la IA generativa\protect\footnotemark}
\end{figure}
\footnotetext{Diapositiva 6.}

\subsection{¿Qué es el aprendizaje profundo?}

Antes de entender el aprendizaje profundo, es necesario aclarar tres niveles de conceptos:

\begin{importantbox}{La jerarquía AI $\supset$ ML $\supset$ DL}
\begin{itemize}
\item \textbf{Artificial Intelligence (AI)}: cualquier técnica que permite a una computadora imitar el comportamiento humano
\item \textbf{Machine Learning (ML)}: un método en el que, en lugar de programar explícitamente, se deja que la máquina aprenda patrones a partir de los datos
\item \textbf{Deep Learning (DL)}: un subconjunto de ML que usa redes neuronales profundas para extraer patrones a partir de datos crudos
\end{itemize}
Idea central: \textbf{enseñar a la computadora a aprender directamente de los datos crudos cómo realizar una tarea} (learn a task directly from raw data).
\end{importantbox}

\begin{figure}[H]
\centering
\includegraphics[width=0.85\textwidth]{slides-images/slide-007.jpg}
\caption{La relación anidada entre AI, Machine Learning y Deep Learning\protect\footnotemark}
\end{figure}
\footnotetext{Diapositiva 7.}

La esencia de \textbf{Intelligence} (inteligencia) es la capacidad de procesar información para guiar decisiones futuras. Artificial Intelligence consiste en construir algoritmos artificiales para llevar a cabo ese mismo proceso: usar datos para impulsar decisiones. Machine Learning no le indica directamente a la computadora cómo procesar los datos, sino que deja que descubra patrones en ellos de forma autónoma. Deep Learning, por su parte, lleva este proceso un paso más allá al usar redes neuronales profundas.

\subsection{Estructura del curso}

MIT 6.S191 dura una semana e incluye los siguientes módulos:

\begin{table}[H]
\centering
\begin{tabular}{lll}
\toprule
\textbf{Fecha} & \textbf{Tema de la clase} & \textbf{Laboratorio} \\
\midrule
Day 1 & Introduction to Deep Learning & Lab 1: Generación de música \\
Day 2 & Deep Sequence Modeling & Lab 2: Visión por computadora \\
Day 3 & Deep Computer Vision & Lab 3: Modelos de lenguaje de gran escala \\
Day 4 & Deep Generative Modeling & -- \\
Day 5 & Deep Reinforcement Learning & Final Project \\
\bottomrule
\end{tabular}
\caption{Panorama del calendario del curso}
\end{table}

Entre los puntos destacados de este año se encuentran: laboratorios compatibles simultáneamente con los marcos de trabajo TensorFlow y PyTorch, así como un laboratorio de LLM completamente nuevo (ajuste fino de un modelo Gemma de 2,000 millones de parámetros y construcción de un evaluador de IA).

\subsection{Resumen de la sección}

El aprendizaje profundo ha logrado avances asombrosos en la última década: de la generación de píxeles borrosos a la clonación de voz en tiempo real y a la conversación dinámica. El curso MIT 6.S191 busca, mediante una semana de entrenamiento intensivo, que los estudiantes dominen las técnicas fundamentales que impulsan estos avances. Todo el curso gira en torno a una idea central: enseñar a la computadora a aprender tareas directamente a partir de datos crudos.
